
\documentclass{beamer}
\usepackage{ctex, hyperref}
\usepackage[T1]{fontenc}

% other packages
\usepackage{latexsym,amsmath,xcolor,multicol,booktabs,calligra}
\usepackage{graphicx,pstricks,listings,stackengine}

\author{{王笙铭}}
\institute{\href{https://www.zstu.edu.cn/}{浙江理工大学信息科学与工程学院}}
\title{ZSTU Beamer Theme}
\subtitle{使用 \LaTeX 制作精致的学术演示文档}
\date{\today}
\usepackage{zstu_beamer}

% defs
\def\cmd#1{\texttt{\color{blue}\footnotesize $\backslash$#1}}   % 反斜杠的命令文本显示为蓝色
\def\env#1{\texttt{\color{red}\footnotesize #1}}
\definecolor{deepblue}{RGB}{0,103,172}
\definecolor{deepred}{rgb}{220,20,60}
\definecolor{deepgreen}{rgb}{0,0.5,0}
\definecolor{halfgray}{gray}{0.55}

\lstset{
    basicstyle=\ttfamily\small,
    keywordstyle=\bfseries\color{deepblue},
    emphstyle=\ttfamily\color{deepblue},   % Custom highlighting style
    stringstyle=\color{deepgreen},
    numbers=left,
    numberstyle=\small\color{halfgray},
    rulesepcolor=\color{red!20!green!20!blue!20},
    frame=shadowbox,
}


\begin{document}

\kaishu
\begin{frame}
    \titlepage
    \begin{figure}[htpb]
        \begin{center}
            \includegraphics[width=0.4\linewidth]{pic/zstu-logo2.png}
        \end{center}
    \end{figure}
\end{frame}

\begin{frame}
    \tableofcontents[sectionstyle=show,subsectionstyle=show/shaded/hide,subsubsectionstyle=show/shaded/hide]
\end{frame}


% 内容从这里开始
\section{课题背景}

\begin{frame}{用Beamer很高大上？}
    \begin{itemize}[<+-| alert@+>] 
    % 除了alert，手动在里面插 \pause 也能达成“一页一单元”的效果
    
        \item 好多学校都有专属的Beamer主题，ZSTU也不能没有！
        \item 中文支持请选择 Xe\LaTeX{} 编译选项
        \item 原始项目的Overleaf地址位于 \url{https://cn.overleaf.com/latex/templates/fzu-beamer-theme-fuzhou-university/bsqkxvgrpfcn}
    \end{itemize}
\end{frame}


\section{研究现状}

\subsection{Beamer主题分类}

\begin{frame}
    \begin{itemize}
        \item \LaTeX{} 模版有很多，可以在这里找到 \newline \url{https://cn.overleaf.com/gallery}
        \item 各高校的Beamer模板最初来源自 \newline \url{https://www.latexstudio.net/archives/4051.html}
    \end{itemize}
\end{frame}


\section{研究内容}

\subsection{美化主题}

\begin{frame}{这一份主题与原始的THU Beamer Theme区别在于}
    \begin{itemize}
        \item 顶栏的小点变成一行而不是多行
        \item 中文采用楷书
        \item 更多该模板的功能可以参考 \url{https://www.latexstudio.net/archives/4051.html}
        \item 如果你想从头开始学习Beamer，非常建议观看这个视频 \url{https://www.bilibili.com/video/BV1Sd4y157cJ/}
        \item 下面我们列举出了一些Beamer的用法，供你参考
    \end{itemize}
\end{frame}

\subsection{使用Beamer来制作演示文档}
\begin{frame}{Why Beamer}
    \begin{itemize}
        \item \LaTeX 广泛用于学术界，期刊会议论文模板
    \end{itemize}
    \begin{table}[h]
        \centering
        \begin{tabular}{c|c}
            Microsoft\textsuperscript{\textregistered}  Word & \LaTeX \\
            \hline
            文字处理工具 & 专业排版软件 \\
            容易上手，简单直观 & 容易上手，进阶困难 \\
            所见即所得 & {\color{deepblue}{所想即所得}}  \\
            高级功能不易掌握 & 进阶难，但一般用不到 \\
            处理长文档需要丰富经验 & 长短文档的处理基本无异 \\
            需要花费大量时间调整格式 & 可以专注于内容本身 \\
            公式排版差强人意 & 尤其擅长公式排版 \\
            二进制格式，兼容性差 & 文本文件，易读、稳定 \\
            付费商业许可 & {\color{deepblue}{开源，自由，免费}} \\
        \end{tabular}
    \end{table}
\end{frame}

\begin{frame}{数学公式插入}
    \begin{itemize}
        \item 由SOGI系统框图，可得如下表达式：
            \[\left\{ \left[ V\left( s \right) -V'\left( s \right) \right] k-V'\left( s \right) 
        \frac{\omega _{\mathrm{n}}}{s} \right\} \frac{\omega _{\mathrm{n}}}{s}=V'\left( s \right) \]
        \item SOGI中$v'(t)$的传递函数为:
            \[H_{\mathrm{d}}\left( s \right) =\frac{V'\left( s \right)}{V\left( s \right)}
            =\frac{k\omega _{\mathrm{n}}s}{s^2+k\omega _{\mathrm{n}}s+{\omega _{\mathrm{n}}}^2}\]
        \item SOGI中虚拟正交信号
            \footnote{脚注测试：用Park变换把合成矢量投影到两个正交直流分量上$u\left( t \right) =K_{\mathrm{I}}\int_0^t{e\left( t \right) dt}$}
        $qv'(t)$的传递函数为：
            \[H_{\mathrm{q}}\left( s \right) =\frac{qV'\left( s \right)}{V\left( s \right)}=\frac{V'\left( s \right)
             \frac{\omega _{\mathrm{n}}}{s}}{V\left( s \right)}=
             \frac{{k\omega _{\mathrm{n}}}^2}{s^2+k\omega _{\mathrm{n}}s+{\omega _{\mathrm{n}}}^2}\]
    \end{itemize}
\end{frame}


\begin{frame}{Align环境(带编号)与图片插入测试}
    \begin{itemize}
        \item Park变换表达式：
    \end{itemize}
    \begin{align}
        \left[\begin{array}{l} \nonumber
            V_d \\
            V_q
            \end{array}\right]=\left[\begin{array}{c}
            V_\alpha \cos (\theta)+V_\beta \sin (\theta) \\
            -V_\alpha \sin (\theta)+V_\beta \cos (\theta)
            \end{array}\right]
            =\left[\begin{array}{cc}   
            \cos (\theta) & \sin (\theta) \\
            -\sin (\theta) & \cos (\theta)
            \end{array}\right]\left[\begin{array}{c}
            V_\alpha \\
            V_\beta
            \end{array}\right]
    \end{align}
    \begin{figure}[htbp]
        \centering
        \includegraphics[width=0.5\textwidth]{pic/Frequency-Response.png}
        \caption{ SOGI在连续时间下的频率响应$H(j\omega)$(Bode Diagram)\label{fig1}} % 后文引用：\ref{}
    \end{figure}

\end{frame}

\begin{frame}
    \begin{exampleblock}{编号多行公式}
        % Taken from Mathmode.tex
        \begin{multline}
            A=\lim_{n\rightarrow\infty}\Delta x\left(a^{2}+\left(a^{2}+2a\Delta x+\left(\Delta x\right)^{2}\right)\right.\label{eq:reset}\\
            +\left(a^{2}+2\cdot2a\Delta x+2^{2}\left(\Delta x\right)^{2}\right)\\
            +\left(a^{2}+2\cdot3a\Delta x+3^{2}\left(\Delta x\right)^{2}\right)\\
            +\ldots\\
            \left.+\left(a^{2}+2\cdot(n-1)a\Delta x+(n-1)^{2}\left(\Delta x\right)^{2}\right)\right)\\
            =\frac{1}{3}\left(b^{3}-a^{3}\right)
        \end{multline}
    \end{exampleblock}
\end{frame}

\begin{frame}{图形与分栏}
    % From thuthesis user guide.
    \begin{minipage}[c]{0.3\linewidth}
        \psset{unit=0.8cm}
        \begin{pspicture}(-1.75,-3)(3.25,4)
            \psline[linewidth=0.25pt](0,0)(0,4)
            \rput[tl]{0}(0.2,2){$\vec e_z$}
            \rput[tr]{0}(-0.9,1.4){$\vec e$}
            \rput[tl]{0}(2.8,-1.1){$\vec C_{ptm{ext}}$}
            \rput[br]{0}(-0.3,2.1){$\theta$}
            \rput{25}(0,0){%
            \psframe[fillstyle=solid,fillcolor=lightgray,linewidth=.8pt](-0.1,-3.2)(0.1,0)}
            \rput{25}(0,0){%
            \psellipse[fillstyle=solid,fillcolor=yellow,linewidth=3pt](0,0)(1.5,0.5)}
            \rput{25}(0,0){%
            \psframe[fillstyle=solid,fillcolor=lightgray,linewidth=.8pt](-0.1,0)(0.1,3.2)}
            \rput{25}(0,0){\psline[linecolor=red,linewidth=1.5pt]{->}(0,0)(0.,2)}
%           \psRotation{0}(0,3.5){$\dot\phi$}
%           \psRotation{25}(-1.2,2.6){$\dot\psi$}
            \psline[linecolor=red,linewidth=1.25pt]{->}(0,0)(0,2)
            \psline[linecolor=red,linewidth=1.25pt]{->}(0,0)(3,-1)
            \psline[linecolor=red,linewidth=1.25pt]{->}(0,0)(2.85,-0.95)
            \psarc{->}{2.1}{90}{112.5}
            \rput[bl](.1,.01){C}
        \end{pspicture}
    \end{minipage}\hspace{1cm}
    \begin{minipage}{0.5\linewidth}
        \medskip
        %\hspace{2cm}
        \begin{figure}[h]
            \centering
            \includegraphics[height=.4\textheight]{pic/dtmf.pdf}
        \end{figure}
    \end{minipage}
\end{frame}

\begin{frame}[fragile]{\LaTeX{} 常用命令}
    \begin{exampleblock}{命令}
        \centering
        \footnotesize
        \begin{tabular}{llll}
            \cmd{chapter} & \cmd{section} & \cmd{subsection} & \cmd{paragraph} \\
            章 & 节 & 小节 & 带题头段落 \\\hline
            \cmd{centering} & \cmd{emph} & \cmd{verb} & \cmd{url} \\
            居中对齐 & 强调 & 原样输出 & 超链接 \\\hline
            \cmd{footnote} & \cmd{item} & \cmd{caption} & \cmd{includegraphics} \\
            脚注 & 列表条目 & 标题 & 插入图片 \\\hline
            \cmd{label} & \cmd{cite} & \cmd{ref} \\
            标号 & 引用参考文献 & 引用图表公式等\\\hline
        \end{tabular}
    \end{exampleblock}
    \begin{exampleblock}{环境}
        \centering
        \footnotesize
        \begin{tabular}{lll}
            \env{table} & \env{figure} & \env{equation}\\
            表格 & 图片 & 公式 \\\hline
            \env{itemize} & \env{enumerate} & \env{description}\\
            无编号列表 & 编号列表 & 描述 \\\hline
        \end{tabular}
    \end{exampleblock}
\end{frame}

\begin{frame}[fragile]{\LaTeX{} 环境命令举例}
    \begin{minipage}{0.5\linewidth}
\begin{lstlisting}[language=TeX]
\begin{itemize}
  \item A \item B
  \item C
  \begin{itemize}
    \item C-1
  \end{itemize}
\end{itemize}
\end{lstlisting}
    \end{minipage}\hspace{1cm}
    \begin{minipage}{0.3\linewidth}
        \begin{itemize}
            \item A
            \item B
            \item C
            \begin{itemize}
                \item C-1
            \end{itemize}
        \end{itemize}
    \end{minipage}
    \medskip
    \pause
    \begin{minipage}{0.5\linewidth}
\begin{lstlisting}[language=TeX]
\begin{enumerate}
  \item 国民 \item 表率
  \item 社会
  \begin{itemize}
    \item[n+e] 栋梁
  \end{itemize}
\end{enumerate}
\end{lstlisting}
    \end{minipage}\hspace{1cm}
    \begin{minipage}{0.3\linewidth}
        \begin{enumerate}
            \item 国民
            \item 表率
            \item 社会
            \begin{itemize}
                \item[\href{https://caihanlin.com/hobbies}{Lance}] 的彩蛋
            \end{itemize}
        \end{enumerate}
    \end{minipage}
\end{frame}

\begin{frame}[fragile]{\LaTeX{} 数学公式}
    \begin{columns}
        \begin{column}{.55\textwidth}
\begin{lstlisting}[language=TeX]
$V = \frac{4}{3}\pi r^3$

\[
  V = \frac{4}{3}\pi r^3
\]

\begin{equation}
  \label{eq:vsphere}
  V = \frac{4}{3}\pi r^3
\end{equation}
\end{lstlisting}
        \end{column}
        \begin{column}{.4\textwidth}
            $V = \frac{4}{3}\pi r^3$
            \[
                V = \frac{4}{3}\pi r^3
            \]
            \begin{equation}
                \label{eq:vsphere}
                V = \frac{4}{3}\pi r^3
            \end{equation}
        \end{column}
    \end{columns}
    \begin{itemize}
        \item 学习更多内容，请看 \href{https://zh.wikipedia.org/wiki/Help:数学公式}{\color{deepblue}{这里}}
    \end{itemize}
\end{frame}

\begin{frame}[fragile]
    \begin{columns}
        \column{.6\textwidth}
\begin{lstlisting}[language=TeX]
    \begin{table}[htbp]
      \caption{编号与含义}
      \label{tab:number}
      \centering
      \begin{tabular}{cl}
        \toprule
        编号 & 含义 \\
        \midrule
        1 & 4.0 \\
        2 & 3.7 \\
        \bottomrule
      \end{tabular}
    \end{table}
    公式~(\ref{eq:vsphere}) 的
    编号与含义请参见
    表~\ref{tab:number}。
\end{lstlisting}
        \column{.4\textwidth}
        \begin{table}[htpb]
            \centering
            \caption{编号与含义}
            \label{tab:number}
            \begin{tabular}{cl}\toprule
                编号 & 含义 \\\midrule
                1 & 4.0\\
                2 & 3.7\\\bottomrule
            \end{tabular}
        \end{table}
        \normalsize 公式~(\ref{eq:vsphere})的编号与含义请参见表~\ref{tab:number}。
    \end{columns}
\end{frame}

\begin{frame}{作图}
    \begin{itemize}
        \item 矢量图 eps, ps, pdf
        \begin{itemize}
            \item METAPOST, pstricks, pgf $\ldots$
            \item Xfig, Dia, Visio, Inkscape $\ldots$
            \item Matlab / Excel 等保存为 pdf
        \end{itemize}
        \item 标量图 png, jpg, tiff $\ldots$
        \begin{itemize}
            \item 提高清晰度，避免发虚
            \item 应尽量避免使用
        \end{itemize}
    \end{itemize}
    \begin{figure}[htpb]
        \centering
        \includegraphics[width=0.3\linewidth]{pic/zstu-logo1.png}
        \caption{图片}
    \end{figure}
\end{frame}

\section{计划进度}
\begin{frame}
    \begin{itemize}
        \item 一月：完成文献调研
        \item 二月：复现并评测各种Beamer主题美观程度
        \item 三、四月：制作ZSTU Beamer主题
        \item 五月：论文撰写
    \end{itemize}
\end{frame}

\section{参考文献}

\begin{frame}{参考链接} % [allowframebreaks]
    \begin{itemize}
        \item 各高校的Beamer模板最初来源自 \newline \url{https://www.latexstudio.net/archives/4051.html}
        \item 原始项目的Overleaf地址位于 \newline \url{https://www.overleaf.com/latex/templates/thu-beamer-theme/vwnqmzndvwyb}
        \item 本项目的GitHub地址位于 \url{https://github.com/GuangLun2000/FZU-latex-template}
    \end{itemize}
    
    % 参考文献引用方式：
    % \bibliography{ref}
    % \bibliographystyle{alpha}
    % 如果参考文献太多的话，可以像下面这样调整字体：
    % \tiny\bibliographystyle{alpha}
\end{frame}


\begin{frame}
    \begin{center}
        {\Huge\calligra Thanks you!}
    \end{center}
\end{frame}

\end{document}


